\section{Open Questions and Next Work}
\label{sec:receipts-questions}

\subsection{Known Gaps}
\label{sec:receipts-questions-gaps}

Four structural gaps in the receipts corpus have been identified and are documented
here as open findings rather than closed facts.

\textbf{Gap 1: Ed25519 signing key custody.} All five M\&A receipts carry
\texttt{validator\_signature} fields. However, the conformance authority generation
receipt marks the Ed25519 signature as \texttt{PENDING}. The private key that would
seal the conformance authority module itself---as distinct from the keys used to sign
the M\&A claim outputs---is not defined in the corpus. The custody protocol for the
signing authority private key, including key generation, storage, rotation, and
revocation, is absent. This gap means that while the M\&A claim receipts are signed,
the infrastructure that manufactures those signatures is not itself signed to the same
standard. \textbf{Classification: FUTURE WORK.}

\textbf{Gap 2: Live event log binding.} The \texttt{RECEIPT\_REGISTRY.md} notes that
``production receipts require live data.'' The five M\&A receipts carry
\texttt{target\_log\_hash} BLAKE3 values, but the actual XES or OCEL event log files
from which those hashes were computed are not present in the receipts directory. Without
the original event logs, independent replication of the conformance measurements is not
possible from within this corpus. The \texttt{rec\_wc\_ar\_002.json} receipt carries a
replication \texttt{PASSED} verdict, but the replication event log itself is not archived
in the repository. \textbf{Classification: FUTURE WORK.}

\textbf{Gap 3: FFI and WASM boundary.} The conformance authority generation receipt
lists three implementation components as \texttt{PENDING}: FFI bindings, receipt
serialization in CBOR format (as opposed to the current JSON/JCS form), and the audit
ledger binding. The JavaScript and WebAssembly boundary between the Rust conformance
authority module and browser-based or Node.js-based consumers is unimplemented.
Without this boundary, the conformance authority cannot be invoked from the web-facing
components of the process-intelligence stack. \textbf{Classification: FUTURE WORK.}

\textbf{Gap 4: Board deck materialization.} The \texttt{ma\_deck\_rendering\_authority\_assessment.md}
defines a seven-step board presentation pipeline and marks all steps as \texttt{VERIFIED}
or \texttt{EXECUTABLE}. However, the final \texttt{ma\_deck\_rendering.md} receipt ledger
file does not exist, and the actual board deck artifacts---the \texttt{.pptx} PowerPoint
file and the \texttt{.xlsx} due diligence workbook---have not been materialized in the
repository. The Tera2 templates and the \texttt{pptx-rs}/\texttt{calamine} rendering
libraries are in place, but the rendering execution trigger has not been pulled.
\textbf{Classification: FUTURE WORK.}

\subsection{Deferred Gates}
\label{sec:receipts-questions-deferred}

Three gates that are logically required by the receipts architecture have been deferred
to future work phases. The first is the \textbf{ADVERSARIAL gate}: the
\texttt{RECEIPT\_REGISTRY.md} defines an \texttt{ADVERSARIAL} receipt type for claims
that have been stress-tested against hostile process variants. No adversarial receipts
have been manufactured. The hostile test cases defined in the
\texttt{adversarial/} directory of the research program have not been run against the
conformance authority to produce adversarial receipts. The gate exists in the registry
but has no passing receipts.

The second deferred gate is the \textbf{STANDARDS gate}: the \texttt{RECEIPT\_REGISTRY.md}
defines a \texttt{STANDARDS} receipt type for anchoring XES, OCEL 2.0, BPMN, ISO, SOC2,
and GDPR compliance maps. These standards compliance assessments exist as research documents
in the \texttt{standards/} directory of the research program, but they have not been
promoted to signed, hashed STANDARDS receipts.

The third deferred gate is the \textbf{CBOR serialization gate}: the current receipt
format uses JSON with JCS canonical serialization. The generation receipt notes CBOR
as a pending serialization format, which would allow receipts to be embedded in
constrained environments (IoT edge nodes, embedded audit hardware) that cannot process
full JSON payloads. The CBOR serialization layer has not been implemented.

\subsection{Research Risks}
\label{sec:receipts-questions-risks}

Three research risks are identified as requiring active monitoring. The first is
\textbf{signature authority drift}: if the Ed25519 private key used to sign the M\&A
claim receipts is lost, rotated without a receipt, or compromised, the entire signed
receipt chain loses its chain-of-custody property. The current corpus does not document
a key rotation or recovery protocol. This risk is rated high because the signing
authority is the trust root of the entire board-admissible claim set.

The second risk is \textbf{event log unavailability}: if the event logs identified by
the BLAKE3 \texttt{target\_log\_hash} values in the five M\&A receipts become unavailable
(deleted, encrypted, or access-restricted), independent replication of the conformance
measurements becomes impossible. The receipts prove that a measurement was made, but
cannot reproduce it without the original log. This risk is rated medium for the current
research context but would be rated high in a live M\&A transaction context.

The third risk is \textbf{algorithm obsolescence}: the conformance authority implements
A* alignment (Adriansyah 2014) and token-based replay (van der Aalst formula). Both
algorithms are well-established but represent a specific point in the conformance
checking literature. If the field moves to alignment algorithms with materially
different fitness semantics, receipts issued under the current algorithm may be
compared against receipts issued under a different algorithm, producing apparent
discrepancies that are artifacts of algorithm change rather than process change.
\textbf{AUTHOR THESIS}: This risk motivates the \texttt{process\_model\_hash}
and algorithm identifier fields in the receipt schema, which allow algorithm
versioning to be tracked alongside hash versioning.

\subsection{Next Receipted Work}
\label{sec:receipts-questions-next}

Three concrete next-work items are identified, each of which must itself produce
a receipt upon completion to enter the research program's evidence chain.

The first next-work item is \textbf{live event log archival}: acquire or manufacture
representative XES/OCEL event logs for each of the five M\&A claim domains, archive
them in the receipts corpus, verify that their BLAKE3 hashes match the
\texttt{target\_log\_hash} values in the existing signed receipts, and emit a
\texttt{LOG\_ARCHIVE} receipt confirming the binding. This item is prerequisite for
the full independent replicability property claimed by the corpus.

The second next-work item is \textbf{FFI and WASM boundary implementation}: implement
the JavaScript/WebAssembly boundary for the conformance authority module, implement
CBOR serialization for receipts, implement the audit ledger binding, and emit a
\texttt{WASM\_BOUNDARY} generation receipt confirming the boundary is tested and sealed.

The third next-work item is \textbf{board deck execution}: trigger the Tera2 template
rendering pipeline to materialize the \texttt{.pptx} board deck and \texttt{.xlsx}
due diligence workbook, then emit the \texttt{ma\_deck\_rendering.md} receipt ledger
file confirming the render succeeded and the output artifacts are hash-bound to the
signed M\&A claim receipts that supplied their data. This item closes the final open
gap in the seven-step board presentation pipeline documented in
\texttt{ma\_deck\_rendering\_authority\_assessment.md}.
